\documentclass{article}
\usepackage[T1]{fontenc}
\usepackage{lmodern}
\usepackage{graphicx}
\usepackage{amsmath,amssymb}
\usepackage[a4paper,margin=2cm]{geometry}
\usepackage[absolute,overlay]{textpos}
\setlength{\TPHorizModule}{1mm}
\setlength{\TPVertModule}{1mm}
\usepackage[indonesian]{babel}
\usepackage{hyperref}
\setlength{\emergencystretch}{2em}
% unit: Halaman 29

\begin{document}

% === Halaman 29 (python/native) ===

• Misalkan panjang kunci m = 10, maka 10 huruf pertama plainteks dienkripsi 

dengan kunci K, setiap huruf ke-i menggunakan kunci ki. 

• Untuk 10 karakter berikutnya, kembali menggunakan pola enkripsi yang sama.

   Contoh: kunci = sony    (m = 4)

 Plainteks: 

thisplaintextissecret

 Kunci: 

sonysonysonysonysonys

   Setiap huruf plainteks dienkripsi dengan huruf kunci di bawahnya dengan  cipher 

abjad tunggal (misalnya menggunakan Caesar Cipher):

                (t+s) mod 26

                  (h + o) mod 26


(i + n) mod 26

                   dst…

\end{document}
